% =====================================================================
% Appendix I (app:related): extended related work (one-column appendix).
% Long version of Sec. 2: every citation cut from the main text (round 2
% also cut there the Sortformer diarization clause, the multilingual
% turn-taking clause and the repeated JSTAR and Seed citations; all are
% covered here), how cstbench eval relates to existing toolkits and
% metrics, the source of every cell of Table 1 (tab:positioning; cells
% checked against the cited papers), and how the paper uses this work,
% incl. the optional Hibiki-Zero baseline and the coupling controls.
% Speech corpora are compared in App. A (tab:corpora); the data paragraph
% below only names them. Planned parts carry a \draftnote.
% =====================================================================
\section{Extended Related Work}\label{app:related}

This appendix extends \cref{sec:related}: it surveys each line of work with its citations, relates \scorer to existing evaluation tools, and gives the source of every cell of \cref{tab:positioning}.

\paragraph{Simultaneous translation policies.}
A simultaneous system decides after each input step whether to read more input or to write output. The fixed wait-$k$ policy trails the source by $k$ words and comes with Average Lagging (AL; \citealp{ma2019stacl}). Learned policies use monotonic attention, from MILk, which also introduced Differentiable Average Lagging \citep{arivazhagan2019monotonic}, to monotonic multihead attention \citep{ma2020monotonic}; wait-$k$ and monotonic multihead attention were adapted to speech input, together with a computation-aware latency metric \citep{ma2020simulmt}. Offline models run simultaneously by local agreement, which commits the prefix on which consecutive outputs agree \citep{liu2020lowlatency,polak2022cuni,machacek2023turning}, or by attention: EDAtt and AlignAtt decide from encoder--decoder attention, or from the alignment derived from it, whether to emit \citep{papi2023attention,papi2023alignatt}, and AlignAtt also drives decoder-only LLM translators \citep{fuxa2026alignatt4llm}. Other policies commit at learned meaning units \citep{zhang2020learning} or learn the segmentation of the speech jointly with translation \citep{zhang2023endtoend}. Re-translation revises the displayed text instead \citep{arivazhagan2020retranslation}; as \task scores append-only output, such systems take part through a commit rule (\cref{sec:exp-setup}).

\paragraph{Streaming systems, LLMs and speech output.}
StreamAtt extends attention-based policies to unbounded audio by selecting which audio and text history to keep, and introduces \laal \citep{papi2024streamatt}; a review of 110 SimulST papers finds that most still use human pre-segmented speech \citep{papi2025how}. Speech LLMs translate unbounded speech as a multi-turn dialogue with a managed KV cache \citep{ouyang2025infinisst} or with a chunkwise causal encoder \citep{ouyang2025cmu}, fold segmentation, policy and translation into one model through speech chain-of-thought \citep{guo2025streamuni}, or act as interpreting agents \citep{cheng2024towards}. StreamSpeech, SeamlessStreaming, Hibiki and Hibiki-Zero also emit speech \citep{zhang2024streamspeech,seamless2023seamless,labiausse2025highfidelity,labiausse2026simultaneous}: Hibiki processes source and target as synchronous streams, trained with per-word delays from contextual alignment \citep{labiausse2025highfidelity}, and Hibiki-Zero dispenses with word-level aligned data by optimising latency with reinforcement learning \citep{labiausse2026simultaneous}. All of these systems are built and evaluated for one direction at a time, and none attributes its output to a speaker.

\paragraph{Evaluating simultaneous and streaming translation.}
SimulEval simulates simultaneous translation of pre-segmented input and reports lagging metrics \citep{ma2020simuleval}; simulstream evaluates long-form streaming translation, also with re-translation \citep{gaido2025simulstream}. AL underestimates the latency of over-generating systems, which Length-Adaptive Average Lagging (LAAL) corrects \citep{papi2022overgeneration}. Stream-level latency resegments the output against the references \citep{iranzosanchez2021stream}, as mwerSegmenter does for quality \citep{matusov2005evaluating}, and \laal applies LAAL to such resegmented output \citep{papi2024streamatt}; a meta-evaluation uncovers a segmentation-related bias of these metrics and introduces YAAL, LongYAAL and the SoftSegmenter resegmenter, implemented in OmniSTEval \citep{polak2025better}. IWSLT moved from sentence-level lagging metrics under a 2-second AL constraint \citep{ahmad2024findings} to unsegmented long-form audio, with \laal as the only latency metric in 2025 \citep{abdulmumin2025findings} and LongYAAL as the primary one in 2026 \citep{adelani2026speech}. All of these tools and metrics score one speaker translated in one direction. \scorer keeps stream-level resegmentation and \laal, but resegments each output language onto the turns of its own direction and adds what a listener in a conversation experiences: the wait after each turn (\eo, which SimulEval reports per pre-segmented input), the delay at each change of direction (switch latency, SimulEval's StartOffset taken at such changes) and output in the wrong direction (\cref{sec:bench-scoring,app:scorer}).
\draftnote{planned in \scorer: switch latency and the wrong-direction detector}

\paragraph{Joint transcription, translation and speaker attribution.}
Streaming ASR and simultaneous ST can share an encoder, with two synchronized decoders in which intermediate ASR results guide the translation policy \citep{chen2021direct}, or one decoder can interleave both, ordered by token-level serialized output training \citep{papi2023token}; StreamSpeech also emits its streaming transcript \citep{zhang2024streamspeech}. Serialized output training (SOT) transcribes overlapping speakers one after another with a single output layer \citep{kanda2020serialized}; its token-level form (t-SOT) emits the tokens of overlapping speakers in time order and marks each switch between virtual output channels with a special token \citep{kanda2022streaming}, and token-level speaker embeddings attribute these tokens while streaming \citep{kanda2022streamingspeaker}. End-to-end speaker-attributed ASR also counts and identifies the speakers \citep{kanda2020joint,kanda2021end}. Speaker embeddings let a streaming transducer ST model detect speaker changes, to drive speaker-aware synthesis of the translation \citep{wang2025streaming}. DiariST builds on t-SOT and token-level speaker embeddings for streaming translation with diarization and scores speaker-attributed BLEU \citep{yang2024diarist}; STAC-ST emits transcripts, translations and speaker-turn tokens in one serialized output, offline \citep{zuluagagomez2023end}; and JSTAR streams joint ASR and ST with a fast--slow cascaded transducer encoder and multi-objective training \citep{moritz2025transcribing}.

\paragraph{Diarization.}
End-to-end neural diarization (EEND) maps frames to multi-label speaker activity with a permutation-free objective and thus handles overlap \citep{fujita2019end,fujita2019endselfattention}; attractors extend it to unknown numbers of speakers \citep{horiguchi2020end}, and a speaker-tracing buffer to streaming input \citep{xue2021online}. Sortformer instead supervises speakers in order of arrival \citep{park2025sortformer}, and Streaming Sortformer keeps the arrival-ordered speaker indices consistent online with a speaker cache \citep{medennikov2025streaming}. Transducers infer speaker roles \citep{elshafey2019joint} or emit speaker-turn tokens that drive streaming diarization \citep{xia2022turn}, while clustering pipelines such as pyannote \citep{bredin2023pyannote} and LLM post-correction of speaker labels \citep{wang2024diarizationlm} work offline. Two-party benchmarks hold monolingual telephone calls \citep{godfrey1992switchboard,cieri2004fisher,canavan1997callhome}, meeting benchmarks include single-device recordings \citep{carletta2006ami,vinnikov2024notsofar}, and simulated conversations add controlled overlap \citep{chen2020continuous,cosentino2020librimix,landini2022from}; none has two speakers who use different languages. Our diarization baselines are pyannote and Streaming Sortformer (\cref{sec:exp-aux}), and the activity head of \model uses arrival-order slots (\cref{sec:model}).

\paragraph{Turn-taking and full-duplex models.}
Conversation analysis describes turns allocated at transition-relevance places with minimal gap and overlap \citep{sacks1974simplest}, yet measured gaps and overlaps are both frequent \citep{heldner2010pauses}; computational turn-taking is reviewed by \citet{skantze2021turntaking}. Voice Activity Projection predicts the upcoming voice activity of both speakers \citep{ekstedt2022voice} and TurnGPT predicts turn shifts from text \citep{ekstedt2020turngpt}. Monolingual turn-taking models transfer poorly to other languages, and one trained on English, Mandarin and Japanese dialogues matches them \citep{inoue2024multilingual}, but every dialogue it learns from is held in one language. Full-duplex dialogue models represent both sides of a conversation on one timeline, as two audio channels \citep{nguyen2023generative}, as parallel audio streams with time-aligned text \citep{defossez2024moshi}, or by running an LLM synchronously with a real-time clock \citep{veluri2024beyond}; delayed-stream and step-synchronous decoders emit text on the audio clock \citep{zeghidour2025streaming,seide2024speech}. \model adopts one clock-synchronous token stream for the transcripts and translations of both speakers (\cref{sec:model}).

\paragraph{Conversational translation systems and data.}
Machine interpreting of conversations is a long-standing goal \citep{waibel1996interactive,wahlster2000verbmobil,lewis2015skype}, and JANUS-II already compared push-to-talk with cross-talk recording \citep{lavie1996translation}. \Cref{app:taxi} checks public bilingual conversational speech corpora against the requirements of a real \task test set (\cref{tab:corpora}): MSLT keeps one side of each bilingual call \citep{federmann2016microsoft,federmann2017microsoft}, DRAL pairs fragments of one-language conversations with re-enactments in the other language \citep{ward2023dialogs,ward2024dialogs,avila2023towards}, InCroMin records meetings mediated by machine translation with one track per participant \citep{cechovic2025corpus}, the Verbmobil~II interpreter volumes need a paid licence \citep{bas2000verbmobil2}, and Fisher--CALLHOME has one source language \citep{post2013improved}. Chat translation pairs an English-speaking agent with a German-speaking customer \citep{farajian2020findings}, or English with Chinese \citep{liang2021modeling}, as \task pairs its speakers, but in text. Code-switched speech translation and data \citep{weller2022end,alastruey2023towards,hamed2022arzen,shankar2025costa,yan2025csfleurs,zhou2025csdialogue} alternate languages within a speaker rather than between speakers, which our scope excludes (\cref{sec:intro}); still, one bidirectional end-to-end model translates code-switched speech well even without code-switched training data \citep{weller2022end}, which supports one model for both directions.

\paragraph{Sources of the cells of \cref{tab:positioning}.}
Each cell follows the cited paper; \pmark{} marks a property that is shown or partly covered but not evaluated, and n.r.\ one that the paper does not report.
\begin{itemize}\setlength{\itemsep}{1pt}
  \item \emph{SimulST/StreamST}: segmented benchmarks such as MuST-C \citep{digangi2019must} and the long-form IWSLT 2025--2026 test sets \citep{abdulmumin2025findings,adelani2026speech} are public, and some systems release weights.
  \item \emph{SeamlessStreaming} translates between many languages into text and speech, is evaluated on FLEURS \citep{conneau2023fleurs} and releases its weights \citep{seamless2023seamless}.
  \item \emph{StreamSpeech} translates into English, also emitting its transcript and speech, is evaluated on CVSS \citep{jia2022cvss} and releases its weights \citep{zhang2024streamspeech}.
  \item \emph{Hibiki} translates French and \emph{Hibiki-Zero} French, Spanish, Portuguese and German into English, in text and speech; both release weights, and Hibiki-Zero also releases a 45-hour evaluation benchmark \citep{labiausse2025highfidelity,labiausse2026simultaneous}.
  \item \emph{STAC-ST} is offline; its two-speaker test channels merge the two sides of Spanish Fisher--CALLHOME calls and contain cross-talk (11--12\% overlap, marked with cross-talk tokens), and it releases scripts but no models \citep{zuluagagomez2023end}.
  \item \emph{DiariST} translates multi-speaker Chinese meetings into English with diarization; it is scored on single-channel audio (the first channel of the far-field array, and a mix of the headset signals) with natural overlap, and releases its evaluation set, DiariST-AliMeeting, with offline baselines and the scorer, but not its streaming model \citep{yang2024diarist}.
  \item \emph{JSTAR} works on the five-channel microphone array of smart glasses and tells the wearer from the partner by the direction of speech; its single output sequence carries wearer and partner tags, and each speaker's speech is translated into the other's language (Spanish--English); it is tested on in-house real conversations and on FLEURS with simulated overlap, and releases neither data nor model \citep{moritz2025transcribing}.
  \item \emph{Seed LiveInterpret~2.0} depicts a two-speaker Chinese--English conversation (its Fig.~2) and is trained for multi-speaker discrimination, but is evaluated one direction at a time, hence \pmark{} for speakers$\times$languages and for the choice of direction; its main test set, RealSI, consists of clips cut from public videos, with public annotations and video links \citep{cheng2024towards}; the paper does not report its input channels, and the system is closed \citep{cheng2025seed}.
  \item \emph{Ours}: TAXI is free under a scientific licence but may not be redistributed, so we release build scripts and checksums that rebuild \taxi from its original distribution \citep{bas2001taxi,reichel2016bas}; \syn is released under CC licences; overlap is rendered (L1, \cref{sec:bench-timing}); and \model is released with its weights.
\end{itemize}
\draftnote{planned in our row: the L1 renders, \syn{} (incl.\ \pair{Ko}{En}), transcript and diarization scoring, and \model{} with its weights}

\paragraph{Use in this paper.}
Several of these systems serve as baselines (\cref{sec:exp-setup}; \cref{app:baselines} lists checkpoints and wrappers). StreamSpeech and, optionally, Hibiki-Zero, whose weights are public (CC BY-NC-SA 4.0) and which also emits text, run as \deen-only text baselines (\cref{tab:conditions}). \scorer reuses stream-level resegmentation and \laal \citep{papi2024streamatt}; \model reuses meaning units \citep{zhang2020learning}, Hibiki's contextual alignment \citep{labiausse2025highfidelity}, the time-ordered serialization of t-SOT \citep{kanda2022streaming} and arrival-order speakers \citep{park2025sortformer} (\cref{sec:model}). As \model decides turns, directions and translations in one stream, whether this coupling helps is tested by a decoupled twin and a same-backbone cascade (\cref{sec:exp-ablation}).
\draftnote{planned: the baseline runs, Hibiki-Zero among them only as an option, the decoupled twin, and \model{} beyond its streaming ASR backbone}
